\begin{resumo}

Este documento apresenta o desenvolvimento do Ratobô, um robô da categoria Micromouse projetado com o objetivo de mapear e resolver 
fisicamente labirintos de dimensões 4x4, 8x8 e 16x16 de maneira autônoma, integrado a uma aplicação web para armazenamento e exibição de telemetria em tempo real. 
A metodologia aplicada engloba a gestão baseada na estrutura analítica do projeto e a divisão técnica nas áreas de eletrônica, software, estrutura e energia. 
O hardware é centralizado em um microcontrolador ESP32, utilizando locomoção diferencial com motores DC N20 equipados com encoders, ponte H MX1508, e três sensores 
ultrassônicos HC-SR04, sendo alimentado por uma bateria de 12V regulada e monitorada pelo módulo INA219. No software, a navegação autônoma emprega a busca em 
profundidade para exploração e o algoritmo Floodfill para o retorno ,enquanto a interface web foi desenvolvida com React, Node.js e PostgreSQL, comunicando-se 
via MQTT e WebSockets. Como resultados esperados a partir da validação planejada, o sistema deve garantir que o robô conclua os percursos respeitando rigorosas restrições 
físicas, permitindo que a aplicação exiba instantaneamente métricas de desempenho como consumo, velocidade e o trajeto em um mapa dinâmico.

\vspace{\onelineskip}
\noindent
\textbf{Palavras-chaves}:  \imprimirpalavrachaveum, \imprimirpalavrachavedois
\end{resumo}
